% Generated draft: requires mathematical and source review.
\documentclass[11pt]{article}
\usepackage[margin=25mm]{geometry}
\usepackage{fontspec}
\setmainfont{Latin Modern Roman}
\usepackage{amsmath,amssymb,mathtools,unicode-math}
\setmathfont{Latin Modern Math}
\usepackage{enumitem,needspace,xurl,hyperref}
\setlength{\parindent}{0pt}
\setlength{\parskip}{5pt}
\setlength{\emergencystretch}{3em}
\newcommand{\N}{\mathbb{N}}
\newcommand{\Z}{\mathbb{Z}}
\newcommand{\Q}{\mathbb{Q}}
\newcommand{\R}{\mathbb{R}}
\newcommand{\C}{\mathbb{C}}
\newcommand{\F}{\mathbb{F}}
\newcommand{\A}{\mathbb{A}}
\newcommand{\PP}{\mathbb{P}}
\newcommand{\E}{\mathbb{E}}
\newcommand{\eps}{\varepsilon}
\newcommand{\dd}{\,\mathrm{d}}
\begin{document}
\section*{UnsolvedMath ID: 800640}
\textbf{Source problem code:} catalogue entry
\subsection{Definitions and mathematical statement}
The supplied record does not specify the ambient class of spaces, the permitted cardinalities of covers, or the definitions of its cover classes. The following gives the usual open-cover interpretation, with these unresolved qualifications made explicit. Let \(\mathscr X\) denote the intended, unspecified class of topological spaces. For \(X\in\mathscr X\), let \(\mathcal O_*(X)\) denote the source's intended class of admissible open covers. Under the usual convention, each \(\mathcal U\in\mathcal O_*(X)\) is a family of open subsets of \(X\) such that \(\bigcup\mathcal U=X\) and \(X\notin\mathcal U\). Whether admissible covers must be countable is unresolved in the supplied record.

Using this interpretation, define the following subclasses of \(\mathcal O_*(X)\):
\begin{itemize}
\item \(\Omega(X)\) consists of the \(\omega\)-covers: for every finite subset \(F\subseteq X\), there is \(U\in\mathcal U\) with \(F\subseteq U\).
\item \(\Gamma(X)\) consists of the \(\gamma\)-covers: \(\mathcal U\) is infinite and, for every \(x\in X\), the family \(\{U\in\mathcal U:x\notin U\}\) is finite.
\item \(\mathrm T(X)\) consists of the \(\tau\)-covers: for every \(x\in X\), the family \(\{U\in\mathcal U:x\in U\}\) is infinite, and for every \(x,y\in X\), at least one of
\[
\{U\in\mathcal U:x\in U,\ y\notin U\},\qquad
\{U\in\mathcal U:y\in U,\ x\notin U\}
\]
is finite.
\end{itemize}
For classes \(\mathcal A(X)\) and \(\mathcal B(X)\) of covers, write
\[
X\models\binom{\mathcal A}{\mathcal B}
\quad\Longleftrightarrow\quad
(\forall\mathcal U\in\mathcal A(X))
(\exists\mathcal V\subseteq\mathcal U)\;
\mathcal V\in\mathcal B(X).
\]

The source question is whether, throughout its intended class of spaces and with its intended cover conventions,
\[
(\forall X\in\mathscr X)\qquad
X\models\binom{\Omega}{\Gamma}
\quad\Longleftrightarrow\quad
X\models\binom{\Omega}{\mathrm T}.
\]
The inclusion \(\Gamma(X)\subseteq\mathrm T(X)\) gives the forward implication. Thus the requested implication is that every space in the intended class for which every admissible \(\omega\)-cover contains a \(\tau\)-subcover also has a \(\gamma\)-subcover of every admissible \(\omega\)-cover. The unspecified ambient class and admissibility conventions must be recovered before this becomes a fully determined assertion.

For context, \(\mathsf S_1(\Omega,\Gamma)\) means that for every sequence \((\mathcal U_n)_{n\in\N}\) of members of \(\Omega(X)\), one can choose \(U_n\in\mathcal U_n\) for each \(n\) so that \(\{U_n:n\in\N\}\in\Gamma(X)\). The supplied source attributes to reference [GN], in its intended setting, the equivalence
\[
\mathsf S_1(\Omega,\Gamma)\Longleftrightarrow\binom{\Omega}{\Gamma}.
\]
It also reports that all examples then considered which failed \(\binom{\Omega}{\Gamma}\) also failed \(\binom{\Omega}{\mathrm T}\). This is historical source context, not a resolution of the question or an independently verified present-day assessment.
\subsection{Short English statement}
In the source's intended class of spaces, does the property that every open \(\omega\)-cover contains a \(\tau\)-subcover imply that every such cover contains a \(\gamma\)-subcover? The converse follows from \(\Gamma\subseteq\mathrm T\). The supplied record leaves the ambient spaces and cover cardinality conventions unspecified. Its label \emph{open} is supplied registry metadata, not an independently verified current status.
\subsection{Sources}
\begin{itemize}
\item Supplied local UnsolvedMath registry record: permanent ID \(800640\), exact source problem code \texttt{catalogue entry}, in the AMR Open Problem Lists collection. The record attributes the problem to Tsaban, with proposed year 2003, and labels it \emph{open}. Its supplied creation and update date is 2026-10-16; upstream revision: \texttt{d371ee917ac9647fbb06cd6fc529deec6f3dda82}.
\item Tsaban, \emph{Selection Principles in Mathematics: A Milestone of Open Problems} (2003), Problem 12. Supplied source URL: \url{https://arxiv.org/abs/math/0312182}. The record reports access on 2026-10-16 and extraction by \texttt{arxiv\_latex}; these are supplied provenance statements, not a claim of independent source access or verification.
\item The title and background retain an attribution/context pointer to Section 4 of the reference labelled \texttt{spm1}. No bibliographic identification of that reference is supplied.
\item The problem text cites reference labels [tau], [tautau], [splittability], [strongdiags], [Hdim], and [MilNonGamma] for studies of \(\tau\)-covers. It attributes \(\mathsf S_1(\Omega,\Gamma)\Longleftrightarrow\binom{\Omega}{\Gamma}\) to [GN]. The supplied record does not identify these references bibliographically.
\end{itemize}
\end{document}
